\subsection{Natural Number Proofs}
To prove $\forall n \in \N. P$, we of course again use induction:

\shade{blue}{Base Case} Show $P[n \mapsto 0]$

\shade{green}{Step Case} Let $m \in \N$ be arbitrary and not free in $P$. We then assume that $P[n \mapsto m]$ and show that $P[n \mapsto m + 1]$

Or the same as a natural deduction rule:
\[
    \begin{prooftree}
        \hypo{\Gamma \vdash P[n \mapsto 0]}
        \hypo{\Gamma, P[n \mapsto m] \vdash P[n \mapsto m + 1]}
        \infer2[$m$ not free in $\Gamma, P$]{\Gamma \vdash \forall n \in \N. P}
    \end{prooftree}
\]


\subsubsection{Induction over the natural numbers}
In Haskell, we can also define all the natural numbers using
\mint{haskell}+data Nat = Zero | Succ Nat deriving (Eq, Ord, Show)+

Thus the natural numbers are (isomorphic to) the set
\[
    \texttt{Nat} = \{ \texttt{Zero}, \texttt{Succ Zero}, \texttt{Succ (Succ Zero)}, \ldots \}
\]

The data type provides two crucial rules for constructing members of \texttt{Nat}:
\begin{itemize}
    \item $\texttt{Zero} \in \texttt{Nat}$
    \item If $x \in \texttt{Nat}$, then $\texttt{Succ} x \in \texttt{Nat}$
\end{itemize}

The induction stated as a natural deduction rule:
\[
    \begin{prooftree}
        \hypo{\Gamma \vdash P[n \mapsto \texttt{Zero}]}
        \hypo{\Gamma, P[n \mapsto m] \vdash P[n \mapsto \texttt{Succ}\ m]}
        \infer2[$m$ not free in $\Gamma, P$]{\Gamma \vdash \forall n \in \texttt{Nat}. P}
    \end{prooftree}
\]


\subsubsection{Lists}
A possible data type for lists in Haskell is:
\mint{haskell}+data L t = Nil | Cons t (L t)+

A natural deduction rule for induction over lists is:
\[
    \begin{prooftree}
        \hypo{\Gamma \vdash P[xs \mapsto \texttt{Nil}]}
        \hypo{\Gamma, P[xs \mapsto ys] \vdash P[xs \mapsto \texttt{Cons}\ y\ ys]}
        \infer2[$y$, $ys$ not free in $\Gamma, P$]{\Gamma \vdash \forall xs \in \texttt{L}\ t. P}
    \end{prooftree}
\]


\subsubsection{Trees}
A possible data type for trees in Haskell is:
\mint{haskell}+data Tree t = Leaf | Node t (Tree t) (Tree t)+

A natural deduction rule for induction over trees is:
\[
    \begin{prooftree}
        \hypo{\Gamma \vdash P[x \mapsto \texttt{Leaf}]}
        \hypo{\Gamma, P[x \mapsto l] \vdash P[xs \mapsto \texttt{Node}\ a\ l\ r]}
        \infer2[$a$, $l$, $r$ not free in $\Gamma, P$]{\Gamma \vdash \forall x \in \texttt{Tree}\ t. P}
    \end{prooftree}
\]


\subsubsection{Structural Induction}
Induction is based on the structure of terms
\mint{haskell}+data T t = Leaf t | Node1 (T t) | Node2 t (T t) (T t)+

\shade{blue}{Base Case} $T_0 = \{ \texttt{Leaf}\ a \divider a \in t \}$

\shade{green}{Step Case} $T_i = T_{i - 1} \cup \{ \texttt{Node1}\ s \divider s \in T_{i - 1} \} \cup \{ \texttt{Node2}\ a\ l\ r \divider a \in t \text{ and } l, r \in T_{i - 1} \}$

A natural deduction rule structural induction is:
\[
    \begin{prooftree}
        \hypo{\Gamma \vdash P[x \mapsto \texttt{Leaf}\ a]}
        \hypo{\Gamma, P[x \mapsto s] \vdash P[xs \mapsto \texttt{Node1}\ s]}
        \hypo{\Gamma, P[x \mapsto l], P[x \mapsto r] \vdash P[ \mapsto \texttt{Node2}\ a\ l\ r]}
        \infer3[$(*)$]{\Gamma \vdash \forall x \in \texttt{T}\ t. P}
    \end{prooftree}
\]
(*) $a$, $l$, $r$, $s$ not free in $\Gamma, P$
